\chapter{Происхождение квадратичной флуктуации \texorpdfstring{$A_{(2)}$}{A(2)}}
\label{ch:v7-r2-generalized-fluctuation-seed-origin-gate}

\section{Вопрос и нулевая гипотеза}

Архитектурная развилка оставила маршрут без новых фермионов: отказаться от
строгого первого порядка и использовать обобщённую внутреннюю флуктуацию.
До вычисления потенциала необходимо проверить, способен ли уже допущенный
родитель сам активировать нелинейный член
\begin{equation}
 A_{(2)}=\sum_{j,k}\widehat a_j a_k[[D,b_k],\widehat b_j].
 \label{eq:v7-a2-seed-source}
\end{equation}
Нулевая гипотеза: текущий оператор поддержан только на стандартных рёбрах
$u,d,e$, проходит первый порядок и потому даёт $A_{(2)}=0$.

\section{Наследование от полного поля Тома VII}

Исправленный носитель имеет вид
\begin{equation}
 \mathcal E_{\rm phys}=E_{\rm aff}\otimes\Lambda_{\rm ch},
 \qquad \Lambda_{\rm ch}=\operatorname{span}_{\mathbb C}\{u,d,e\}.
 \label{eq:v7-a2-admitted-carrier}
\end{equation}
Семейный множитель коммутирует с конечной алгеброй. Поэтому
\begin{equation}
 [[X_a\otimes T_a,\pi(b)],\pi^o(c)]
 =X_a\otimes[[T_a,\pi(b)],\pi^o(c)]=0.
 \label{eq:v7-a2-family-factorization}
\end{equation}
Линейность даёт
\begin{equation}
 D_{\rm adm}=\sum_{a=u,d,e}X_a\otimes T_a
 \quad\Longrightarrow\quad A_{(2)}[D_{\rm adm}]=0.
 \label{eq:v7-a2-admitted-zero}
\end{equation}
Ранги $24\to21$ и Hodge-неустойчивость меняют коэффициенты и ранг на
прежней опоре, но не создают скрытого нарушения первого порядка.

\section{Центральный аудит и положительный контроль}

На одном представителе каждого неприводимого бимодуля введены
\begin{equation}
 P_s^L|_{\mathcal H_{ij}}=\delta_{s,i},
 \qquad P_t^R|_{\mathcal H_{ij}}=\delta_{t,j}.
 \label{eq:v7-a2-central-projectors}
\end{equation}
Для всех девяти пар центральных идемпотентов получено
\begin{equation}
 \max_{s,t}\left\|[[D_{\rm adm},P_s^L],P_t^R]\right\|=0.
 \label{eq:v7-a2-admitted-residual}
\end{equation}
Положительный контроль использует запрещённые диагонали $L_L-u_R$ и
$Q_L-e_R$ и даёт
\begin{equation}
 \max_{s,t}\left\|[[D_{R_2},P_s^L],P_t^R]\right\|>0.
 \label{eq:v7-a2-forbidden-residual}
\end{equation}

\section{Одна унитарная обобщённая флуктуация}

Для $L=\pi(u)$, $R=\pi^o(u)$ и $U=LR$:
\begin{align}
 D^u&=UDU^*,\nonumber\\
 A_{(1)}&=LDL^*-D,\qquad
 \widetilde A_{(1)}=RDR^*-D,\nonumber\\
 A_{(2)}&=D^u-D-A_{(1)}-\widetilde A_{(1)}.
 \label{eq:v7-a2-unitary-decomposition}
\end{align}
Машинная проверка даёт
\begin{equation}
 \|A_{(2)}[D_{\rm adm}]\|<10^{-15},
 \label{eq:v7-a2-unitary-admitted-zero}
\end{equation}
тогда как для $D_{R_2}$ норма ненулевая. При добавлении
$\varepsilon D_{R_2}$ она линейна по $|\varepsilon|$ с остатком ниже
$10^{-15}$.

\section{Сверка с прежней ветвью}

В Pati--Salam ненулевой $A_{(2)}$ возникал после явного расширения алгебры
и конечного seed. Нормированный полугрупповой аудит затем доказывал
составные тождества, но слишком общий seed заполнял весь 72-мерный
симметричный канал. Общая теория показывает, что $A_{(2)}$ исчезает при
первом порядке \cite{v7a2-generalized}; Pati--Salam является новой
геометрией, а не флуктуацией неизменённой Стандартной модели
\cite{v7a2-pati-salam}.

\begin{theorem}[Нулевой $A_{(2)}$ допущенного родителя]
На носителе $E_{\rm aff}\otimes\Lambda_{\rm ch}$ любой оператор,
поддержанный на рёбрах $u,d,e$, удовлетворяет первому порядку до семейного
тензорного множителя. Поэтому квадратичная часть любой обобщённой внутренней
флуктуации равна нулю.
\end{theorem}

\begin{nogocriterion}[Круговой seed]
Если ненулевой $A_{(2)}$ появляется только после включения в исходный $D$
запрещённого кварк-лептонного блока, конструкция не выводит $R_2$, а
предполагает новый сектор входными данными.
\end{nogocriterion}

\begin{protocol}
\begin{enumerate}
 \item допущенная опора: \textbf{$u,d,e$};
 \item максимальный первопорядковый остаток: \textbf{ноль};
 \item $A_{(2)}$ от допущенного родителя: \textbf{ноль};
 \item семейный подъём меняет вывод: \textbf{нет};
 \item запрещённый $R_2$-seed активирует $A_{(2)}$: \textbf{да};
 \item $R_2$ получен без вставки нового сектора: \textbf{нет};
 \item обобщённый маршрут как вывод: \textbf{закрыт круговостью};
 \item следующий гейт: \textbf{Real- и аномальное завершение зеркального
 цикла}.
\end{enumerate}
\end{protocol}

Машинный сертификат сохранён в
\path{s2t_v7_r2_generalized_fluctuation_seed_origin_gate_results.json}.

\begin{thebibliography}{9}
\bibitem{v7a2-generalized}
A.~H.~Chamseddine, A.~Connes, W.~D.~van Suijlekom,
\emph{Inner Fluctuations in Noncommutative Geometry without the First Order Condition},
Journal of Geometry and Physics 73 (2013), 222--234; arXiv:1304.7583.
\bibitem{v7a2-pati-salam}
A.~H.~Chamseddine, A.~Connes, W.~D.~van Suijlekom,
\emph{Beyond the Spectral Standard Model: Emergence of Pati--Salam Unification},
Journal of High Energy Physics 2013, 132; arXiv:1304.8050.
\end{thebibliography}